\documentclass[UTF8,a4paper,11pt,openany,fontset=fandol]{ctexbook}
\usepackage{xr}
\usepackage[a4paper,margin=24mm,headheight=15pt]{geometry}
\usepackage{amsmath,amssymb,bm,mathtools,booktabs,longtable,tabularx,array,graphicx}
\usepackage{enumitem,xcolor,fancyhdr,needspace}
\usepackage[numbers,sort&compress]{natbib}
\usepackage{hyperref,bookmark}
\usepackage{tikz}
\usetikzlibrary{arrows.meta,positioning}
\setCJKmainfont{FandolSong-Regular.otf}[BoldFont=FandolSong-Bold.otf,ItalicFont=FandolKai-Regular.otf]
\setCJKsansfont{FandolHei-Regular.otf}[BoldFont=FandolHei-Bold.otf]
\setCJKmonofont{FandolFang-Regular.otf}
\setmainfont{texgyretermes-regular.otf}[BoldFont=texgyretermes-bold.otf,ItalicFont=texgyretermes-italic.otf,BoldItalicFont=texgyretermes-bolditalic.otf]\setsansfont{texgyreheros-regular.otf}[BoldFont=texgyreheros-bold.otf]
\definecolor{ink}{HTML}{193D55}
\ctexset{chapter={format=\sffamily\huge\bfseries\color{ink},beforeskip=0pt,afterskip=20pt},section={format=\sffamily\Large\bfseries\color{ink}},subsection={format=\sffamily\large\bfseries},subsubsection={format=\sffamily\normalsize\bfseries}}
\hypersetup{colorlinks=true,linkcolor=ink,citecolor=ink,urlcolor=ink,bookmarksnumbered=true,bookmarksdepth=2,pdftitle={Hall MHD与双流体磁重联：模型推导、数值方法与文献调研}}
\setlength{\emergencystretch}{4em}\setlength{\parskip}{3pt}\linespread{1.12}\raggedbottom
\setlist{nosep,leftmargin=2em}\setcounter{tocdepth}{1}\setcounter{secnumdepth}{2}
\pagestyle{fancy}\fancyhf{}\fancyhead[L]{\small\nouppercase{\leftmark}}\fancyfoot[C]{\thepage}
\newcommand{\vct}[1]{\bm{#1}}\newcommand{\dd}{\mathrm d}\newcommand{\VA}{V_{\!A}}\newcommand{\J}{\vct J}\newcommand{\B}{\vct B}\newcommand{\E}{\vct E}\newcommand{\vel}{\vct u}\newcommand{\grad}{\nabla}\newcommand{\R}{\mathcal R}
\newcommand{\note}[1]{\par\smallskip\noindent\textsf{\textbf{学习提示：}}#1\par\smallskip}
\newcommand{\assume}[1]{\par\smallskip\noindent\textsf{\textbf{假设与边界：}}#1\par\smallskip}
\newcolumntype{P}[1]{>{\raggedright\arraybackslash}p{#1}}
\numberwithin{equation}{chapter}

\let\cleardoublepage\clearpage

\tracinglostchars=3

\externaldocument[V1-][nocite]{01_theory_numerics}[01_theory_numerics.pdf]
\externaldocument[V2-][nocite]{02_paper_results}[02_paper_results.pdf]
\hypersetup{pdftitle={Hall MHD与双流体磁重联：已解决与未解决的问题}}
\begin{document}\frontmatter\begin{titlepage}\vspace*{22mm}
{\sffamily\Large Hall MHD 与双流体磁重联调研\par}\vspace{14mm}
{\sffamily\Huge\bfseries\color{ink}已解决与未解决的问题\par}\vspace{10mm}
{\sffamily\large 三卷本 · 第3卷\par}\vfill
分卷更新：2026年11月12日。\\[5pt]
全文核读证据：139篇，核读更新至2026年11月12日。\\[5pt]
本卷逐项列出已解决结果与未解决问题，标明作者、年份、稳定文献键及可点击引文；原文缺口与本报告综合判断分别标注。后半卷保留更详细的模型边界、争议与可检验研究设计。
\end{titlepage}\chapter*{本卷阅读与查阅方式}
本卷逐项列出已解决结果与未解决问题，标明作者、年份、稳定文献键及可点击引文；原文缺口与本报告综合判断分别标注。后半卷保留更详细的模型边界、争议与可检验研究设计。

三卷分别编排目录和参考文献，文献键保持一致；同一论文在不同卷的引文编号可能不同。标明“基础卷、结果卷、问题卷”的交叉引用会跳转到对应PDF；各卷应保存在同一目录。若阅读器限制跨文件跳转，可按卷名和章节/公式号查阅。全部原文位于项目根目录的papers文件夹。

总索引：\href{run:../paper_index.html}{打开可搜索文献索引}。
\par\href{run:01_theory_numerics.pdf}{打开《理论与数值基础》}。
\par\href{run:02_paper_results.pdf}{打开《不同文章的结果》}。
\tableofcontents\mainmatter
\chapter{实验依据、适用条件与已解决和未解决的问题}\label{ch:summary}
本章先汇总实验依据与模型适用条件，再把每项已解决结果和未决问题对应到具体论文。完整方程和算法见《理论与数值基础》，逐篇核读档案见《不同文章的结果》。后半卷进一步展开模型边界与研究设计。
\section{实验看到了什么：哪些支持Hall物理，哪些仍未验证}
下面汇总已核读原文。实验出现Hall特征，通常验证的是某一物理环节；不能自动认证整套流体闭合及所有尺度的预测。
\begin{longtable}{P{27mm}P{62mm}P{59mm}}
\toprule 实验/文献 & 实际现象及其支持的环节 & 尚不能据此认定什么\\\midrule\endfirsthead
\toprule 实验/文献 & 实际现象及其支持的环节 & 尚不能据此认定什么\\\midrule\endhead
MRX四极场，Ren2005 & 探针测得面外四极磁场，横向尺度与$d_i$同阶；支持解耦电流系统 & 未闭合X点全部欧姆预算，不能证明Hall项直接断线。\\
MRX双尺度层，L119 & 由磁场及密度反演的电子通道窄于离子通道；电子层约$8d_e$，与部分二维计算有差距 & 尺度趋势吻合不等于绝对层宽、排出速度与压力支撑已解释。\\
MRX导引场，GuideExp2012、L193 & 导引场增强伴随率降低、面内结构改变及下游磁压堆积 & 局地Hall效应与装置边界背压未被独立隔离。\\
MRX离子增能，Yoo2013 & 面内电势、离子流与光谱联测；先加速、下游流速趋饱和后温升 & 电功不等于不可逆热化率；碰撞、再磁化及波粒机制仍需区分。\\
MRX碰撞磁岛，E448 & 亚离子电阻层存在Hall四极与低$S$磁岛链 & 不能将无碰撞或单流体的临界阈值直接搬用；具体触发机制仍待检验。\\
MAST合并，Stanier2013 & Hall模拟结合环形压缩及真实诊断弦，可解释密度双峰形态与先后演化 & 峰值和时间尺度未逐时定量复现；等温比和增强耗散限制热分配判断。\\
\bottomrule
\end{longtable}
\citep{Ren2005,L119,GuideExp2012,L193,Yoo2013,E448,Stanier2013}

还应保留能提醒模型扩展的实验：MRX中性密度增加时Hall形态可近似保留而离子喷流变慢，需要加入中性摩擦；MAIZE的四极密度目前只支持待进一步验证的Hall解释，配套三维主计算是电阻MHD；MMS冷/热离子可在同一处分处不同磁化状态，需要更细的物种划分。这些都限制了“见到某图样就认定单离子Hall模型足够”的推断。\citep{L190,F703,L227}

\section{Hall MHD的适用条件怎样判断}
应针对目标区域及目标诊断检查，而不是只给整套装置贴上“Hall”标签。
\begin{enumerate}
\item \textbf{低频、准中性。}目标尺度通常满足$L\gg\lambda_D$，不以电荷分离和电子等离子振荡为研究对象；忽略位移电流还要检查相应低速电磁排序。鞘层、光波或高频电荷分离不属于标准Hall模型的描述范围。
\item \textbf{Hall修正对目标过程重要。}直接检查$|\J|/(enU)$及欧姆预算；在典型Alfvén排序下约为$d_i/L$。宏观盒子很大不妨碍局地薄层进入Hall区。强导场时声回旋半径与有限离子回旋尺度也可能控制解耦，不能仅用$L=d_i$作普适边界。
\item \textbf{被删的电子惯性确实较小。}检查$(m_e/e)D_e\vel_e$相对实际保留项的量级；$d_e/\delta$是常用提示，但不能替代频率、流速和几何检查。若电子扩散区由惯性主导，应采用相应惯性模型。
\item \textbf{压力和热流闭合足以回答所问问题。}标量Hall适合部分离子尺度场和平均流结构；若目标是非对角压力、强各向异性、束流或速度空间热化，需张量/动力学信息或经验证的闭合。
\item \textbf{物种、碰撞和边界正确。}Hall MHD不必严格无碰撞；电子仍可有电阻，同时离子、电子出现解耦。多离子、中性拖曳、辐射电离或装置背压重要时，须把这些过程纳入，而不是靠调大电阻统包。
\item \textbf{电子破冻结的处理可核查。}研究重联必须明确由电阻、电子惯性、压力张量或其他已说明项支撑。仅有无耗散正压Hall输运不足以提供一般的电子断线机制。
\end{enumerate}
因此“Hall MHD可用”应表述为：在指定尺度、闭合与边界下，可用来预测哪些量；它不是整场模拟从离子尺度到电子热化都有效的保证。《理论与数值基础》第1章提供删项估计，《理论与数值基础》第2章提供数值诊断，《不同文章的结果》第2章提供实验测量细节。

\section{已解决到哪一步：逐项对应原始论文}
\label{sec:guide-solved}
以下每项给出作者引文、年份和稳定文献键。点击引文可进入参考文献，再由“全文核读档案”进入《不同文章的结果》第4章；用文献键搜索HTML索引可直接打开本地原文及图号记录。
这里的“已解决”指在明确条件下建立了可检验机制或获得了可靠结果，不表示所有几何和参数下均已解决。

\begin{longtable}{P{29mm}P{63mm}P{56mm}}
\toprule 问题 & 结果与具体论文 & 证据边界\\\midrule\endfirsthead
\toprule 问题 & 结果与具体论文 & 证据边界\\\midrule\endhead
为何必须区分离子与电子流 &
\citet{Ren2005}（2005，Ren2005）图2、4测到MRX四极磁场及碰撞性关联，支持解耦电流系统；\citet{Birn2001}（2001，Birn2001）提供相关跨模型基准。 &
Ren2005的电子流及Hall电场含反演/估计，未闭合X点欧姆预算；不能写成Hall项直接断线已获验证。\\
已启动重联能否得到快演化 &
\citet{Birn2001}（2001，Birn2001）图1中Hall、混合与PIC通量曲线接近，建立了有限种子二维非线性阶段的基准。 &
强种子刻意跳过线性触发；电子层和喷流宽度仍不同。该文不证明自然启动或普适0.1。\\
同一参数下为何可有快慢两支 &
\citet{Cassak2005}（2005，Cassak2005）图3、4给出不同历史及迟滞；\citet{Cassak2010}（2010，Cassak2010）进一步在关闭电子惯性的Hall算例中保留双稳态。 &
依赖具体片层、边界和耗散；2010现象学模型未自洽选定层长，不能给出全部环境的触发阈值。\\
MAST双峰密度怎样形成 &
\citet{Stanier2013}（2013，Stanier2013）将Hall密度响应、环形内侧压缩和实验诊断弦结合，解释双峰及外峰随合并消失的顺序。 &
是形态和顺序的解释；峰值、时钟未逐时定量匹配，且采用固定温比和增强耗散。\\
X线为何单向或双向扩展 &
\citet{Shepherd2012}（2012，Shepherd2012）图3--6显示弱导场主要向电子流动方向扩展，强导场转为双向近导场Alfvén传播。 &
只检验了电子承担初始全部电流等条件；阈值和速度有诊断误差，后期斜模破坏简单图像。\\
流体能否改善电子层描述 &
\citet{Wang2015}（2015，Wang2015）图4、8--10显示十矩可改善部分局地电子流和压力结构；\citet{F611}（2016，F611）图1--7用各向异性流体电子再现长电流层。 &
Wang2015未重现完整时史及所有压力峰值；F611是混合动理学离子模型。这些不是标准标量Hall已完全解决电子层的证据。\\
\bottomrule
\end{longtable}

例如第一行的实验验证、第二行的基准一致和最后一行的闭合改善，属于不同类型的进展。不能把它们叠加成“某一套Hall代码已经同时完成全部验证”。

\section{未解决什么：原文缺口与本报告判断分别标注}
\label{sec:guide-unsolved}
以下问题均对应已核读论文。本节所说“未解决”限于这些论文提供的证据，不能据此断言所有后续文献都没有答案。“原文缺口”指文中明确保留或结果直接显示的不足；“本报告综合判断”是跨论文比较后提出的研究问题，不冒充作者原话。

\subsection{电子扩散区的绝对层宽、速度和欧姆预算}
\textbf{对应文章：}\citet{L119}（2008，L119，
\emph{Experimental study of the Hall effect and electron diffusion region during magnetic reconnection in a laboratory plasma}）；
\citet{Wang2015}（2015，Wang2015）。

\textbf{原文缺口：}L119图13的电场预算没有测全压力张量与涨落项；图16--18的电子层较所比较二维模型更宽、排出更慢，经验摩擦估计仍有差距。
Wang2015图4的重联时史不同；图8--10在等磁通、不同时间比较时，部分非对角压力峰值及下游欧姆项仍不匹配。
\textbf{本报告综合判断：}同为电子尺度的趋势一致，尚未构成对层厚、速度和全部非理想项的同时验证；下一步应在相同诊断定义和仪器平均下比较完整预算。

\subsection{同一压力/热流闭合能否跨算例迁移}
\textbf{对应文章：}\citet{Ng2020}（2020，Ng2020，
\emph{An improved ten-moment closure for reconnection and instabilities}）；
\citet{F611}（2016，F611）；
\citet{Closure2026}（2026，Closure2026）。

\textbf{原文缺口：}Ng2020图1--7及附录A显示，梯度闭合可改善岛合并和部分不稳定性，电子各向异性及跨场热流仍有误差；F611的Harris实现需要二维磁通标签及参考密度修正；Closure2026恢复了部分三维不稳定性，却仍有增长时序、各向异性位置与范围的差异。
\textbf{本报告综合判断：}尚未建立一套无需逐算例调整、又能同时预测率、层结构、热流和三维模态的通用闭合。由某一算例改善不能推出跨几何适用。

\subsection{局地快率能否预测全局完成时间和尺寸效应}
\textbf{对应文章：}\citet{Birn2001}（2001，Birn2001）；
\citet{L241}（2017，L241，
\emph{The role of guide field in magnetic reconnection driven by island coalescence}）；
\citet{Stanier2013}（2013，Stanier2013）。

\textbf{原文证据与缺口：}Birn2001针对已启动单片；L241图1--10中Hall与PIC的岛合并速率、压力阻力和尺寸依赖在弱导场下显著不同，强导场时某些率接近但拓扑和层宽仍不同；Stanier2013中磁通堆积、反弹及岛是否排出受装置几何控制。
\textbf{本报告综合判断：}局地快通量处理能力不足以唯一给出整个磁岛系统完成合并的时间。需要相同初边值、统一上游归一化和累计磁通诊断；这些论文没有共同建立跨系统的普适完成时间定律。

\subsection{有限三维X线的受抑长度是否普适，Hall模型能否独立复现}
\textbf{对应文章：}\citet{F659}（2019，F659，
\emph{Three-Dimensional Magnetic Reconnection With a Spatially Confined X-Line Extent}）；
\citet{F792}（2020，F792，
\emph{Scaling of Magnetic Reconnection With a Limited X-Line Extent}）。

\textbf{原文缺口：}F659图2--9及F792图2--5给出活跃/受抑区与整体率下降；两篇都是PIC研究，电子非对角压力仍支撑破冻结。约$10d_i$受抑尺度来自特定反平行序列，拟合参数、短层瞬态、导场及$\beta$推广仍有限制。
\textbf{本报告综合判断：}尚不能把该尺度写成标准Hall MHD的普适定律。有限长度与压力闭合是否相互改变，是由这两篇及F611共同引出的研究设计，并非三篇已经证实的结论。\citep{F611}

\subsection{离子加速怎样变成不可逆热化，能量怎样分给两物种}
\textbf{对应文章：}\citet{Yoo2013}（2013，Yoo2013，
\emph{Observation of Ion Acceleration and Heating during Collisionless Magnetic Reconnection in a Laboratory Plasma}）；
\citet{L272}（2021，L272）；
\citet{E512}（2011，E512）。

\textbf{原文缺口：}Yoo2013图2--4确立电势、加速与下游升温的空间联系，却保留碰撞、再磁化、波粒及边界等解释。L272图1--7提出电场梯度失磁化途径，但给定场能量展宽和PIC高温关联不能单独唯一确认所有条件下的热化机制；E512在强导场、低$\beta$排序下建立相混合与弱碰撞熵产链，范围又不同。
\textbf{11月12日专题补充：}双温Hall并非完全不能研究加热。\citet{L223}（2016）图3已给出两物种不同热点，\citet{Kaptanoglu2020}以三维双温Hall比较HIT-SI实验，\citet{TenBarge2019MMS}（2019）图1已输出两物种热能预算；但分别受超电阻/输运闭合、未取得的勘误及压力热流闭合限制。\citet{Tanabe2015MAST}的热点位置为前一类模拟提供实验约束，尚未唯一验证局地热化机制。

\textbf{比例必须区分：}\citet{F674}正式版图7的约2:1是离子/电子电场沉积率；同图热能变化约1:1，且含储存与边界焓流。\citet{Yoo2017Heating}非对称MRX沉积更接近；\citet{Yoon2023PhaseSpace}仍明确留下精确离子能量份额和有序/随机增能占比。\textbf{本报告综合判断：}不能用Hall总功为零否定物种增能，也不能用局地正电功替代不可逆热化。条件性两温预算已经可算；普适不可逆加热比仍需动理学信息或经独立检验的闭合，并分别排查模型缺失、数值误差、输运边界和诊断不足。独立专题源与PDF位于项目根目录\texttt{specific problems/离子加热问题/}。

\subsection{自然触发阈值与三维后期状态能否由同一模型预测}
\textbf{对应文章：}\citet{Birn2001}（2001，Birn2001）；
\citet{Cassak2005}（2005，Cassak2005）；
\citet{Cassak2010}（2010，Cassak2010）；
\citet{Shepherd2012}（2012，Shepherd2012）。

\textbf{原文缺口：}GEM基准跳过线性触发并明确保留三维问题；双稳态研究的触发与耗散、片层历史相关，简化模型未闭合层长和变化的上游场；Shepherd2012只在斜模明显相互作用前拟合传播，后期不再满足孤立X线图像。
\textbf{本报告综合判断：}这些研究各解决一个阶段，尚不能直接拼成从任意初始电流片到自然启动、三维扩展及饱和的无参数预测。《不同文章的结果》第1章的启动线与第三维线应联合检验，而不能把阶段性公式无条件串接。



\chapter{模型边界、未解决问题与可检验的研究方向}
\section{先区分方程缺失与尚未算够}
“没有解决”至少有三种含义：方程根本没有这个自由度；方程包含相关项，但闭合没有经过验证；原则上能够研究，但可用计算的尺度、维数和统计时间不足。这三种情况需要不同的下一步。加密网格不能恢复被方程删去的非热分布；更复杂闭合也不能替代真正的三维尺寸扫描。

\begin{longtable}{P{35mm}P{57mm}P{56mm}}
\toprule 问题 & 当前模型的限制 & 可以如何继续\\\midrule
电子非热尾、束流、速度空间多峰 & 标量／有限矩没有唯一分布函数；相同低阶矩可对应不同尾部 & 动理学对照或明确的重构假设；不从温度直接推断加速谱\\
电子扩散区非对角压力 & 普通Hall MHD缺失；十矩大小仍由热流闭合决定 & 同初值比较局地、梯度、非局地闭合，并做欧姆定律残差\\
粒子熵产生与不可逆热化 & 理想流体能量转移不等于微观熵增；等温模型有隐含热库 & 分物种能量预算；碰撞、相混合和数值耗散分别记账\\
全系统触发及等待时间 & 大种子GEM绕过线性触发；边界可能持续供能或阻碍流出 & 噪声、驱动幅度、片厚和边界的独立扫描\\
电子层长度选择 & 局地标度常仍含未决定的层长和上游场 & 同时测长度、局地场、闭合／耗散及全局率\\
三维持续湍流与X线传播 & 二维排除了某些不稳定性；有限第三维限制可用波数 & 第三维长度与分辨率分别收敛；记录前沿和寿命统计\\\bottomrule
\end{longtable}

\section{闭合能否跨算例使用，是目前最具体的难题}
局地松弛可在一个算例调到接近PIC，却在另一系统产生错误的磁岛、各向异性或增长时间。岛合并的证据已显示：改变离子闭合可以改善厚度与压力支撑，但未必同时修复平均率的尺寸标度。\citep{Ng2015,Ng2017,Ng2020}

2026年的三维结果则更明确：恢复低混合漂移和弯曲模，与扩大到不合理区域的各向异性同时发生。\citep{Closure2026}因此有价值的目标不是“再选一个使速率合适的 $k_0$”，而是提前固定闭合，在Harris片、岛合并、不对称片、导引场与三维扰动中共同检验。应把闭合的参数训练算例与验证算例分开。

本报告建议的可判别指标是：增长率与模数、电子／离子压力张量误差、X点欧姆定律分量、混合宽度、峰值和时间平均重联率，以及闭合所产生的能量通量。只有速率接近而压力错误，不能称为模型等价。

\section{“普适0.1”究竟是哪一个量的普适性}
GEM、非对称Hall和岛合并使用的初值、上游定义与统计窗不同。无显式耗散惯性模型的有限Hall扫描也得到速率随参数变化。\citep{Birn2001,Cassak2008,Ng2015,Andres2016}有待澄清的不是一个孤立小数，而是：哪些局地几何或压力约束能选择开口角；这些约束是否对外部驱动、系统尺寸与闭合保持稳定。

可行的Hall／双流体研究应同时给出有量纲电场、全局初始参数归一化和局地上游归一化；先固定 $d_i,d_e$ 与闭合，再改变盒子尺寸，避免改变 $L/d_i$ 时也无意改变分辨率或耗散尺度。还应分别检验层长是输出，还是理论中预设的参数。

\section{快慢分支、磁岛循环和三维湍流是否属于同一相图}
小系统中的快慢双稳态、大系统二维中的循环转换、三维中的持续湍流，目前不能仅凭名称合并为一个完整相图。\citep{Cassak2005,Cassak2010,Huang2011,Huang2024}它们的电阻、超扩散、电子惯性、边界、压力闭合和扫描路径不同。

建议选取统一方程，至少以 $S,L/d_i,B_g/B_{\rm rec}$ 为独立轴，并对每一点从快、慢两种历史启动。先在二维定位寿命与转换，再用相同面内网格加入第三维；观察长期单X线、多磁岛循环和持续湍流是否存在重叠参数区。这是两种流体方法原则上能够推进的问题，但对真实电子层物理的推断仍需动理学对照。

\section{强导引场的低频约化何时失效}
强导引场并不允许无限制地把一切简化成两场。约化模型需满足频率、各向异性、密度扰动和磁化尺度排序；有限电子温度时，压力张量同阶修正也可能不可省。Schep等的原文直接指出，相关传播尺度上的Landau共振会导致流体模型失效；后来的细丝化研究也明确承认无限变薄后的失效。\citep{Schep1994,L049}

因此对观察到的细丝、爆发或饱和，应额外记录最小尺度相对于 $d_e,r_{Le}$ 及沿场自由程的位置，比较流体与动理学的能量分配。单靠总能量守恒和精细图像不足以判定近似仍成立。

\section{实验比较应包含仪器和装置几何}
实验中的导引场堆积、流出受限、驱动电路和有限探针空间分辨率，都会改变同一个局地层的可观测信号。MRX导引场实验已显示局地Hall作用与全局边界作用难以分离。\citep{GuideExp2012}有意义的下一步是在模拟中生成与实验同位置、同时间平均、同空间卷积的合成诊断，而不是把未滤波的二维切片与实验单条剖面直接对照。

\section{三维前沿与湍动的启动竞争}
Hall与双流体已经能研究重联传播，但未统一预测非均匀电流片、有限自由磁能和斜向模相互作用下的前沿。可检验的问题是：在保持局地片厚与扰动不变时，独立改变可用磁通、边界和长宽比，X线伸长转成仅平移的条件是否仍由磁能耗尽决定？这需要能量收支和多组参数扫描，而非单幅传播图。\citep{L275}

十矩三维研究还显示，电子密度噪声从$10^{-10}$改到$10^{-2}$，LHDI启动时间可从约$30$变到$4\Omega_i^{-1}$。磁种子驱动的重联与噪声驱动的不稳定性谁先形成，能改变后续湍动与多X线结构。下一步须同时扫描噪声谱、幅度、闭合与网格；闭合在固定噪声下“吻合一次”，并不足以证明对自然启动的预测能力。\citep{E552}

\section{把来源信息引入流体：新模型应怎样接受检验}
两个电子群的研究把“寻找一个普遍热流公式”改成“增加有限数量的可演化来源标签”。已有低$\beta$固定离子证据说明它可能有效，但真正的研究问题是：何时分群、何时合并、怎样保持总质量动量能量和压力正定、怎样避免标签梯度导致虚假自由膨胀，以及再循环后每群内部多束化如何监测。可用三阶矩重建残差、不同标签选择、不同群数及长期磁通演化共同检验；不能只报某个时间的最佳误差。\citep{F635}

递归Hall撕裂也需要闭环验证：同时测每级$L_*,a_*,B_*,\rho_*$和欧姆平衡，在物理耗散固定下细化，并比较静态本征预测与真实片形成史。当前文献已有“次级增长加快”的证据，却没有在同一非线性数据集上验证Hall修正的临界长宽比。强Hall下次级岛缺失仍需区分物理抑制、寿命过短与漏分辨。\citep{L236,L262}

\section{由强导引场新对照明确下来的四项检验}
\paragraph{外层长度由谁选定。}无快速色散波的快重联反例已存在，但匹配率仍含$B_h,w_v$。下一步可固定$d_e,\rho_s$、耗散和分辨率，分别改变流入盒宽、流出盒长、初始片厚、开放／周期边界；逐时测量$w_v$，而非仅拟合峰率。若某种压力或喷流不稳定性选择了长度，应能预测扫描中何时失效。这属于两流体方法本身可推进的问题。\citep{L216,L217}

\paragraph{大岛双重加速能否通过动理学检验。}已有陀螺流体／动理学非线性对照集中在小$\Delta'$，不能代替大$\Delta'$验证。需要采用相同热尺度、边界、种子及足够高次谐波，分别比较第一轮局地电势增强、两轮之间停滞、第二轮排气开口和最终磁通率；再扫描电子回旋半径、Landau响应及微层分辨率。单独重现最后0.1的峰值会漏掉最有区分力的过程。\citep{F672,E475}

\paragraph{热流闭合能否同时预测时间延迟与能量。}非局地六场模型改善温度平滑，却未复现参考动理学中的加热延迟。可以先在冷离子相同几何比较高阶闭合及更完整动理学，再启用热离子六场，测$Q_\parallel$、温度峰值、磁能释放到热响应的时滞和释放磁能归一化的数值误差。有限矩仍不能唯一决定非热尾部；“自由能被闭合移除”也不能直接给出真实粒子熵增。\citep{L244}

\paragraph{三维磁通旁路与实验全局压缩。}局域体力算例显示磁通绕开层，MRX则显示导引场被装置压缩；两者都说明局地片无法独立代表整体系统。先用同样驱动力、相同磁通库存及同样诊断积分域做二维／三维对照，再在装置几何中加入实际边界和驱动。可检验磁通被重联、被绕流带走和暂时堆积的各自比例。Hall／双流体可以研究这些几何效应，电子层加热则仍需更高矩或动理学检验。\citep{L076,F717,L193}


\section{百篇精读后优先推进的六个问题}
下面的“未解决”均指所读证据尚未给出跨模型、跨参数的充分答案，不声称全世界没有其他研究。每项同时说明Hall/双流体能够推进到哪一步。

\begin{longtable}{P{30mm}P{57mm}P{61mm}}
\toprule 问题 & 已有证据与剩余缺口 & 可执行的下一步\\\midrule\endhead
电子层的真实厚度与慢排出 & MRX约$8d_e$、$0.11V_{\!Ae}$；经验摩擦改善但未闭合压力/涨落预算\citep{L119} & 同装置边界、同探针卷积，比较五矩、十矩、非局地热流；同时匹配宽度、速度、$E$残差，不能只调一条总率曲线\\
三维斜模如何成为非线性X线 & 惯性force-free与电阻Harris线性结果不同；倾角误差不足以区分两选择准则\citep{L224,L268,L166} & 固定同一平衡，扫描电子非理想项与允许波数，比较完整三维模式竞争；区分相速度和前沿速度\\
爆发加速何时被岛/涡旋截断 & 冷模型增加横向分辨率后出现岛；EMHD岛内射流另有KH类失稳\citep{L191,F697,E459} & 幅度定义统一，网格与滤波双扫描；诊断局域X几何、岛出现时间、能量守恒，勿只拟合一个幂指数\\
传播记忆能否在真实磁尾保持 & 非均匀片厚使前沿依赖起始片厚；当前冷电子小质量比模型有限\citep{L294} & 在同一几何加入电子压强/导场和不同载流份额，检验历史效应是否仍存在；装置测量加入有限时间/空间分辨\\
高阶矩能否自洽预测粒子分布 & 最大熵重建正定但21矩仍缺多峰；重建尚非时演验证\citep{E553} & 先检验可实现性和守恒，再从相同初值演化；同时比较压力、热流、分布距离与重联率，而非只拟合已知PIC矩\\
磁能转移如何变成不可逆热 & Hamiltonian能量转移无正定熵源；KREHM给条件性相混合热化机制\citep{E512} & 独立扫描空间与Hermite分辨率、碰撞频率和观察时间；区分数值耗散、压力功和碰撞熵增\\
\bottomrule
\end{longtable}

\subsection{哪些边界不能靠更细Hall网格消除}
标量压强Hall模型没有任意电子非旋转对称压力、俘获/反向粒子群、速度空间相混合和粒子高能尾的自由度。
增加网格能减少离散误差，却不会使方程长出这些未知量。
十矩能提供压力张量通道，但热流仍需闭合；最大熵矩族能增加分布形状的表达能力，却有可实现性和计算代价问题。
KREHM保留电子平行动理学，但其低$\beta$强导场排序又排除了一般反平行高$\beta$情形。
因此一个模型不能仅因在GEM总磁通上成功，就被认为同时解决了磁尾分布、实验加热和高导场湍流。

\subsection{几项值得研究但不能过早宣称的推断}
\citet{F697}中的太阳耀斑时间应用依赖很强的参数外推，部分数值接近相对论或不满足强导场序；需先确认排序，再讨论观测意义。
\citet{L313}关于指数拉伸和电子惯性的框架也需给定具体三维流及边界，才有可验证的时间预测。
\citet{F644}的无显式电子耗散Hall湍流应先量出离散电场残差，才能判断小片重联是否在收敛极限存在。
这些不是否定研究方向，而是给后续模拟设置明确的成功标准。


\Needspace{10\baselineskip}
\section{新增证据留下的具体可研究问题}
\begin{longtable}{P{36mm}P{67mm}P{46mm}}
\toprule 问题 & 现在证据到哪里 & 合适的下一项计算\\\midrule\endfirsthead
\toprule 问题 & 现在证据到哪里 & 合适的下一项计算\\\midrule\endhead
全局率是否趋于平台 & 固定宏观形状时率随尺寸下降，较大Hall域下降变缓；平均区间存在稿件矛盾\citep{F600} & 固定片厚与固定纵横比两套扫描，分别报告局地/全局率、通量阶段和重复扰动误差\\
Hall与PIC何时足够相近 & 强导引场合并中宏观率接近，压力分解和次级岛仍不同\citep{L241} & 用十矩及非局地闭合扫描导引场、离子温度和尺寸，不在每例重新调闭合参数\\
磁岛何时控制净能量转换 & 多种有岛/无岛结果可有相近率；数值耗散与质量比会改变岛\citep{F606,F788} & 同步收敛拓扑、率、压力应变及边界通量；比较形成时间与被排走时间\\
非对称电子层怎样局域化 & 不同非理想项导致不同层长与岛序列\citep{F625} & 匹配电子尺度厚度，再比较张量闭合与超电阻，并单独扫描域长\\
电子单独重联的二次生成 & 大尺度喷流与电子剪切两条路径有二维时序支持\citep{F655} & 多格/$d_e$与滤波收敛、鞍点轨迹、三维及电子温度自由度\\
Full-F爆发与暂态的关系 & 非正规性与非线性谱展宽均出现，但因果量化不足\citep{L314} & 用物理自由能范数算最优增益，控制初始扰动并测非线性转折\\
Hall发电机符号由何控制 & 多螺旋后续事件出现竞争，样本有限\citep{F798} & 保留实际压力剖面和壁面预算，扩大事件数并按前态分类\\
多组分与有限装置 & 冷离子改变分离线电场，中性拖曳降速而层宽未扩\citep{L227,L190} & 多离子/中性独立动量模型，匹配实际边界和碰撞尺度\\
\bottomrule
\end{longtable}

这些问题中，速率、层长、压力输运与大尺度耦合可用Hall或扩展流体继续推进；电子非热尾、微观不可逆加热、尘埃充电和完整多束分布则需额外方程或动力学对照。模型能生成某种形态，不等于其所有微观解释已包含在模型里。

\section{原文一致性审计：保留疑点，不替作者猜参数}
\label{sec:remaining-audit}
本表针对实际核读版本。它记录需要读者注意的公式或标注问题，不表示已证明作者程序错误；版本不同也可能已修正。报告推导均从明确符号和基本方程重建。
\begin{longtable}{P{24mm}P{76mm}P{49mm}}
\toprule 文献/位置 & 原文问题 & 本报告处理\\\midrule\endfirsthead
\toprule 文献/位置 & 原文问题 & 本报告处理\\\midrule\endhead
F600 v1页3--4 & 图1平均磁通阴影约0.5--3，正文写0.05--0.3 & 引用尺寸趋势，精确平均统计待澄清\\
F589 v1页7 & 张角写为$L_i/d_m$，与前文定义及随后公式的倒数关系不符 & 使用$d_m/L_i$，明确为几何重建\\
F606 v1式4 & 若$e>0$，电子压力与惯性项符号应与电子动量相符，印刷式有疑点 & 用本章电子欧姆定律统一符号\\
L314 v2页6 & 写$k=1$同时给$\Delta'=14.3108$，与$2(1/k-k)$不一致 & 不将该组文字参数作为可复现输入\\
L314 v2图11--12 & 坐标写$\hat d_e^2$，图注写$\hat d_e$且行列描述反置 & 保留参数映射不确定性；不据此精确比较跨排序区域\\
E473 v1图6 & 正文与图注左右对应相反 & 依据图内标签：左Full-F，右Full-F+LWL\\
E473/L314能量式 & 磁能写在种类求和中可能重复；个别系数/符号不一致 & 从正定极化二次型及物理磁能重建，不直接抄表\\
F655 v1页6、附录A & 电子/离子Alfvén速度比写反；主事件时刻320.7与360.7混用 & 采用$V_{Ae}/V_{Ai}=\sqrt{m_i/m_e}$；具体数据配对待核验\\
L283 v2页2 & $d_{i0}$式呈现为频率量纲，扰动坐标与$xz$几何不一致 & 用$c/\omega_{pi}$定义尺度，不擅自补齐复现初值\\
F703 v1正文/附录A & 干涉零点基准视线的文字对应不一致 & 保留诊断系统误差，不把投影唯一反演为三维Hall结构\\
\bottomrule
\end{longtable}

正式版F788对噪声、有限电流片寿命和全局几何的讨论比早期稿更谨慎，故本报告以正式版为主要依据。L314引言对F788“欠分辨”的方向性转述与原研究不符，本报告采用F788自己的算例和讨论。这个交叉核对说明，综述不能仅沿引用链复制结论。


\section{沿七条问题线确定下一步，而非按新奇图像选题}
《不同文章的结果》第1章的相反结果经常来自不同问题或不同条件，不能一律当作尚未解决的逻辑矛盾。下表把已有证据与可区分的下一步连接起来；“未解决”指本报告已核读证据仍不足，并非断言世界上没有后续研究。
\begin{longtable}{P{25mm}P{55mm}P{66mm}}
\toprule 问题线 & 当前证据留下的具体缺口 & 有区分力的下一步\\\midrule\endfirsthead
\toprule 问题线 & 当前证据留下的具体缺口 & 有区分力的下一步\\\midrule\endhead
快率（结果卷~\ref{V2-story:fast}） &
色散关联已有反例；层长与供磁仍会随参数调整。 &
保持外层条件可比，分别关闭/改变Hall、惯性和压力响应，同时检查实际X点破冻结预算。\\
全局（结果卷~\ref{V2-story:global}） &
局地快层能与反弹、旁路和尺寸减速共存。 &
固定宏观形状与初始片宽的两种扫描分别进行；统一累计磁通、上游量和完成时间定义。\\
启动（结果卷~\ref{V2-story:onset}） &
次级岛可能促进、阻塞或仅伴随重联；递归临界标度未完全统一。 &
测量每代电流片的局地尺度、寿命与扰动，扫描网格和噪声，区分线性增长与净能量释放。\\
闭合（结果卷~\ref{V2-story:closure}） &
长电子层、热流和三维模态尚不能由一套闭合普遍复现。 &
把F611状态方程与梯度十矩在同一初值比较；检验上游标签、层长、压力散度及能量一致性。\\
第三维（结果卷~\ref{V2-story:three}） &
有限长X线的受抑尺度主要来自PIC特定序列，Hall/双流体能复现到哪一步未定。 &
共同初边值下扫描质量比、导场、$\beta$和边界；同时验证$B_z$输运、活跃长度、压力项与连接变化。\\
热化（结果卷~\ref{V2-story:energy}） &
轨道展宽与高温区关联尚未唯一识别不可逆机制。 &
沿粒子轨道与流体体积双重记账；比较中心矩、压力应变、热流、碰撞熵产和诊断平均。\\
观测（结果卷~\ref{V2-story:diagnostic}） &
嵌套四极和下游升温均有不止一种解释。 &
联合场、流、压力与时间尺度；用多点重建或合成轨道传播坐标及采样误差。\\
\bottomrule
\end{longtable}

本次最直接可由Hall/双流体继续尝试的问题，是\textbf{有限长度抑制与压力闭合是否相互改变}：电子磁通输运能否正确给出活跃/受抑区，而闭合又是否正确给出变薄时间和长层？这是一项由《不同文章的结果》第\ref{V2-story:closure}与《不同文章的结果》第\ref{V2-story:three}节共同引出的研究设计，不是已证实的新规律。若只把PIC压力项换成任意电阻后碰巧得到相近总率，还不足以完成这个检验。\citep{F611,F659,F792}

原书目还缺少若干连接环节，已补登记Ohia2012、Ohia2015和Jain2009，但尚未取得并完成全文核读。它们暂不作为本报告结论证据。原有优先队列中的EMHD撕裂渐近/数值对照、三维电子独有重联、KH跨模型比较和新$\beta$标度也继续保留；不把“列入下一批”写成已经发现或证实其题名所述机制。

\backmatter\bibliographystyle{unsrtnat}\bibliography{03_solved_open}\end{document}
